\section{De actor-critic a métodos puramente basados en valor}

Las clases anteriores presentaron el método actor-critic: mantiene simultáneamente una política (actor) y una función de valor (critic). Q-Learning va más lejos: \textbf{elimina por completo la política explícita} y solo aprende la función Q, de la cual se deriva implícitamente la política.

\begin{knowledgebox}{Basado en valor vs. actor-critic}
\begin{itemize}
    \item \textbf{Actor-Critic}: aprende simultáneamente la política $\pi_\theta$ y la función de valor $\hat{V}$ o $\hat{Q}$.
    \item \textbf{Basado en valor (Q-Learning)}: solo aprende la función Q óptima $Q^*$, y la política queda implícita como $\pi(s) = \arg\max_a Q^*(s, a)$.
\end{itemize}
\end{knowledgebox}

\subsection{Resumen de la sección}

Q-Learning es un método puramente basado en valor que no necesita una red de política explícita.

